
\subsubsection{4.1 Hardware and Software}

All experiments ran on a system with an NVIDIA H100 80GB GPU. Disk capacity was 3.4TB at 100\% utilization during experiments, necessitating checkpoint deletion after each run. Consequently, only final-checkpoint evaluation (step 500) was possible; intermediate checkpoints were not retained. Wall-clock time for the 22 new runs in the h-e1-v2 experiment was approximately 68 minutes total on the H100.

Software: Python 3.10, PyTorch 2.8.0+cu128, HuggingFace Transformers, lm-evaluation-harness. A \texttt{CUDA\textbackslash \_VISIBLE\textbackslash \_DEVICES=0} environment variable fix was required to enable CUDA access in subprocess calls for evaluation.

\subsubsection{4.2 Preliminary Experiment (h-e1)}

Prior to the main experiment, a preliminary experiment (h-e1) was conducted using smaller proxy models (approximately 7M and 16M parameters) with 200 training steps on 5,000 FineWeb documents. This experiment yielded ANOVA p=1.0 and $\eta^{2}\approx 0$ --- no detectable interaction signal at that scale. MMLU 4-shot scores were at the floor (approximately 0.05, consistent with random or below-random for the 10-class format used). Mock data contamination issues were discovered in three successive rounds, requiring re-runs with verified real FineWeb data. These issues are documented and do not affect the h-e1-v2 experiment.

The scope was subsequently reduced to 14M/31M model labels (7.9M/18.1M actual parameters) with 500 training steps and 50,000 source documents.
